{\small
\begin{threeparttable}
\begin{tabular}{llrrr}
\toprule
Group & Occupation & Jobs & Sectors & Index \\
\midrule
High (G1) & Office clerk & 1,674,145 & 87 & 0.938 \\
 & Administrative assistant & 981,313 & 87 & 0.958 \\
 & Production line feeder & 870,424 & 86 & 0.935 \\
 & Public area cleaner & 462,258 & 87 & 0.886 \\
 & Receptionist & 434,935 & 87 & 0.914 \\
\addlinespace[3pt]
Medium (G2) & Construction laborer & 896,163 & 87 & 0.747 \\
 & Cleaner & 765,792 & 87 & 0.841 \\
 & Truck driver & 737,286 & 87 & 0.816 \\
 & Cook & 456,847 & 86 & 0.577 \\
 & Bricklayer & 415,528 & 87 & 0.695 \\
\addlinespace[3pt]
Low (G3) & Retail salesperson & 2,052,934 & 87 & 0.423 \\
 & Cashier & 649,676 & 84 & 0.336 \\
 & Security guard & 528,877 & 86 & 0.267 \\
 & Building doorman & 440,116 & 85 & 0.572 \\
 & General farmworker & 370,982 & 82 & 0.152 \\
\bottomrule
\end{tabular}
\begin{tablenotes}[flushleft]\footnotesize
\item \textit{Note:} Five largest occupations in each group in 2010, by number of private jobs active on December 31. ``Sectors'' is the number of CNAE divisions in which the occupation has at least one job; the index is that of \hyperref[eq:portabilidade]{Equation~(\ref*{eq:portabilidade})}. Source: RAIS, authors' calculations.
\end{tablenotes}
\end{threeparttable}
}
